\documentclass[../../main]{subfiles}

\begin{document}

\section{Ideais e Estruturas Quociente em Anéis}
\label{sec:matematica:algebra:estruturas-algebricas:ideais-quocientes}

Nesta seção consideramos anéis como estruturas sobre a assinatura
\[
	\Sigma = \{+, \cdot : 2\},
\]
e introduzimos subestruturas compatíveis com suas operações, permitindo a construção de novos anéis a partir de relações de equivalência.

\subsection{Ideais}

Seja $(A,+,\cdot)$ um anel. Um subconjunto $I \subseteq A$ é chamado de \textbf{ideal} se satisfaz:

\begin{itemize}
	\item $(I,+)$ é um subgrupo de $(A,+)$, isto é:
	      \begin{itemize}
		      \item $0 \in I$;
		      \item se $a,b \in I$, então $a-b \in I$;
	      \end{itemize}

	\item absorção pela multiplicação:
	      \[
		      \forall a \in A,\ \forall x \in I,\quad a \cdot x \in I \ \text{e}\ x \cdot a \in I.
	      \]
\end{itemize}

Um ideal captura a noção de “subestrutura compatível com a assinatura do anel”.

\subsection{Ideais e congruências}

Todo ideal $I \subseteq A$ induz uma relação de equivalência em $A$ definida por:
\[
	a \sim b \iff a - b \in I.
\]

Essa relação é compatível com as operações do anel, isto é:
\begin{itemize}
	\item se $a \sim b$ e $c \sim d$, então $a+c \sim b+d$;
	\item se $a \sim b$ e $c \sim d$, então $a \cdot c \sim b \cdot d$.
\end{itemize}

\subsection{Anel quociente}

Dado um ideal $I \subseteq A$, o conjunto das classes de equivalência é denotado por:
\[
	A/I = \{[a] : a \in A\}.
\]

Definem-se as operações:
\[
	[a] + [b] = [a+b], \quad [a]\cdot[b] = [a \cdot b].
\]

Essas operações são bem definidas devido à compatibilidade do ideal com a estrutura do anel.

Assim, $(A/I, +, \cdot)$ é novamente um anel.

\subsection{Interpretação estrutural}

O anel quociente $A/I$ pode ser interpretado como a estrutura obtida ao identificar elementos de $A$ que diferem por elementos do ideal $I$.

Isso permite “colapsar” partes do anel preservando sua estrutura algébrica.

\subsection{Ideais e preservação estrutural}

Ideais podem ser vistos como o núcleo de homomorfismos de anéis. Essa relação será formalizada posteriormente, quando introduzirmos morfismos entre estruturas algébricas.

\subsection{Observação conceitual}

A construção de anéis quociente mostra que estruturas algébricas não são apenas conjuntos com operações, mas também objetos que podem ser transformados de forma controlada, preservando sua assinatura e axiomas.

Essa ideia será central para toda a teoria de homomorfismos e para a construção de novos objetos algébricos a partir de estruturas existentes.

\end{document}
